\subsection{\texorpdfstring{THE HOMOMORPHISM \(\operatorname{Hom}_{\mathbb{C}}(\mathrm{E}_1, \mathrm{F}_1) \otimes_{\mathbb{C}} \operatorname{Hom}_{\mathbb{C}}(\mathrm{E}_2, \mathrm{F}_2) \to \operatorname{Hom}_{\mathbb{C}}(\mathrm{E}_1 \otimes_{\mathbb{C}} \mathrm{E}_2, \mathrm{F}_1 \otimes_{\mathbb{C}} \mathrm{F}_2)\)}{The homomorphism from Hom over C of E1,F1 tensor Hom over C of E2,F2 to Hom over C of the corresponding tensor products}}

Let C be a commutative ring and \(E_{1}\) , \(E_{2}\) , \(F_{1}\) , \(F_{2}\) four C-modules; in §3, no. 5, formula (13) we defined a canonical C-module homomorphism

\[
\lambda : \operatorname{Hom} (\mathbf {E} _ {1}, \mathbf {F} _ {1}) \otimes \operatorname{Hom} (\mathbf {E} _ {2}, \mathbf {F} _ {2}) \to \operatorname{Hom} (\mathbf {E} _ {1} \otimes \mathbf {E} _ {2}, \mathbf {F} _ {1} \otimes \mathbf {F} _ {2}). \tag {21}
\]

PROPOSITION 4. When one of the ordered pairs \((\mathbf{E}_1, \mathbf{E}_2)\) , \((\mathbf{E}_1, \mathbf{F}_1)\) , \((\mathbf{E}_2, \mathbf{F}_2)\) consists of finitely generated projective C-modules, the canonical homomorphism (21) is bijective.

It is obviously sufficient to perform the proof for the ordered pairs \((\mathbf{E}_{1}, \mathbf{F}_{1})\) and \((\mathbf{E}_{1}, \mathbf{E}_{2})\) .

We consider first the case of the ordered pair \((\mathbf{E}_1, \mathbf{F}_1)\) ; we fix \(\mathbf{E}_2, \mathbf{F}_1, \mathbf{F}_2\) and write for every C-module \(\mathrm{T}(\mathbf{E}) = \mathrm{Hom}(\mathbf{E}, \mathbf{F}_1) \otimes_{\mathbb{C}} \mathrm{Hom}(\mathbf{E}_2, \mathbf{F}_2)\) and \(\mathrm{T}'(\mathbf{E}) = \mathrm{Hom}(\mathbf{E} \otimes \mathbf{E}_2, \mathbf{F}_1 \otimes \mathbf{F}_2)\) and, for every C-homomorphism \(v: \mathbf{E} \to \mathbf{E}'\) ,

\[
\mathrm{T} (v) = \operatorname{Hom} (v, \mathrm{l} _ {\mathrm{F} _ {1}}) \otimes \mathrm{l} _ {\operatorname{Hom} (\mathrm{E} _ {2}, \mathrm{F} _ {2})}
\]

and

\[
\mathrm{T} ^ {\prime} (v) = \operatorname{Hom} (v \otimes 1 _ {\mathbb {E} _ {2}}, 1 _ {\mathbf {F} _ {1} \otimes \mathbf {F} _ {2}}).
\]

Then Lemmas 4 and 5 (no. 2) (where v is replaced by λ) are valid and are proved by completely analogous methods.

We next fix \(\mathbf{E}_2\) and \(\mathbf{F}_2\) and this time write, for every C-module \(\mathbf{F}\) , \(\mathrm{T}(\mathbf{F}) = \mathrm{Hom}(\mathbf{C},\mathbf{F})\otimes_{\mathbb{C}}\mathrm{Hom}(\mathbf{E}_2,\mathbf{F}_2)\) and \(\mathrm{T}'(\mathrm{F}) = \mathrm{Hom}(\mathbf{C}\otimes \mathbf{E}_2,\mathbf{F}\otimes \mathbf{F}_2)\) and, for every C-homomorphism \(u:\mathrm{F}\to \mathrm{F}'\) ,

\[
\mathrm{T} (u) = \operatorname{Hom} (\mathrm{l} _ {\mathbf {c}}, u) \otimes \mathrm{l} _ {\operatorname{Hom} (\mathbb {E} _ {2}, \mathrm{F} _ {2})}
\]

and

\[
\mathrm{T} ^ {\prime} (u) = \operatorname{Hom} (\mathrm{l} _ {\mathrm{c}} \otimes \mathrm{l} _ {\mathrm{E} _ {2}}, u \otimes \mathrm{l} _ {\mathrm{F} _ {2}}).
\]

This time it is immediately verified that Lemmas 1 and 2 (no. 2) (where \(\lambda\) always replaces \(\nu\) ) are valid.

This being so, we show the proposition first when \(E_{1} = C\) and \(F_{1}\) is projective and finitely generated. The argument of no. 2 (which rests on Lemmas 1 and 2), together with the above remarks, reduces this to proving the proposition when also \(F_{1} = C\) ; then \(\operatorname{Hom}(E_{1}, F_{1})\) , \(E_{1} \otimes E_{2}\) and \(F_{1} \otimes F_{2}\) are identified with C, \(E_{2}\) and \(F_{2}\) respectively (§ 3, no. 4, Proposition 4); the two sides of (21) are then both canonically identified with \(\operatorname{Hom}(\mathbf{E}_{2},\mathbf{F}_{2})\) and, after these identifications, it is verified that \(\lambda\) becomes the identity.

We now suppose that \(F_{1}\) is projective and finitely generated; the argument of no. 2 (depending this time on Lemmas 4 and 5) reduces the proof for \(E_{1}\) any finitely generated projective module to the case where \(E_{1} = C\) , that is to the first case dealt with.

For the ordered pair \((\mathbf{E}_{1}, \mathbf{E}_{2})\) , the procedure is similar, this time applying Lemmas 4 and 5 twice; we leave the details to the reader.

Note that when \(\mathbf{E}_1 = \mathbf{C}^{(I)}\) , \(\mathbf{E}_2 = \mathbf{C}^{(J)}\) are free (finitely generated or not), then \(\operatorname{Hom}(\mathbf{E}_1, \mathbf{F}_1) = \mathbf{F}_1^{\mathrm{I}}\) , \(\operatorname{Hom}(\mathbf{E}_2, \mathbf{F}_2 = \mathbf{F}_2^{\mathrm{J}}\) and

\[
\operatorname{Hom} \left(\mathrm{E} _ {1} \otimes \mathrm{E} _ {2}, \mathrm{F} _ {1} \otimes \mathrm{F} _ {2}\right) = \left(\mathrm{F} _ {1} \otimes \mathrm{F} _ {2}\right) ^ {\mathrm{I} \times \mathrm{J}}
\]

to within canonical isomorphisms and (21) is then identical with a special case of the canonical homomorphism (22) of § 3, no. 7.

When \(E_{2} = C\) , the canonical homomorphism (21) gives, after identifying \(\operatorname{Hom}(E_{2}, F_{2})\) with \(F_{2}\) and \(E_{1} \otimes E_{2}\) with \(E_{1}\) , a canonical homomorphism

\[
\operatorname{Hom} (\mathrm{E}, \mathrm{F}) \otimes \mathrm{G} \rightarrow \operatorname{Hom} (\mathrm{E}, \mathrm{F} \otimes \mathrm{G})\tag{22}
\]

for any three C-modules E, F, G which is just the homomorphism (7) of no. 2 for \(\mathbf{A} = \mathbf{B} = \mathbf{C}\) .

Note that when \(\mathbf{F} = \mathbf{C}\) the canonical homomorphism (22) again gives (11) (no. 2) for the case of a commutative ring.

Suppose now that \(\mathbf{F}_1 = \mathbf{F}_2 = \mathbf{C}\) ; as \(\mathbf{F}_1 \otimes \mathbf{F}_2\) is identified with \(\mathbf{C}\) , there is this time a canonical homomorphism

\[
\mu : \mathbf {E} ^ {*} \otimes \mathbf {F} ^ {*} \rightarrow (\mathbf {E} \otimes \mathbf {F}) ^ {*}\tag{23}
\]

for two C-modules E, F; for \(x^{*} \in \mathbf{E}^{*}\) , \(y^{*} \in \mathbf{F}^{*}\) , the image of \(x^{*} \otimes y^{*}\) under the canonical homomorphism (23) is the linear form \(u\) on \(\mathbf{E} \otimes \mathbf{F}\) such that

\[
u (x \otimes y) = \langle x, x ^ {*} \rangle \langle y, y ^ {*} \rangle .\tag{24}
\]

Moreover, if \(E_{1}\) , \(E_{2}\) , \(F_{1}\) , \(F_{2}\) are four C-modules, \(f: E_{1} \rightarrow E_{2}\) , \(g: F_{1} \rightarrow F_{2}\) two linear mappings, it follows immediately from (24) that the diagram

\[
\begin{array}{c} \mathrm{E} _ {2} ^ {*} \otimes \mathrm{F} _ {2} ^ {*} \xrightarrow {\mu} (\mathrm{E} _ {2} \otimes \mathrm{F} _ {2}) ^ {*} \\ t _ {f} \otimes t _ {g} \Biggl \downarrow \qquad \qquad \qquad \qquad \qquad \qquad \qquad \qquad \qquad \qquad \qquad \qquad \qquad \qquad \qquad \qquad \qquad \qquad \qquad \qquad \qquad \qquad \qquad \qquad \qquad \qquad \qquad \qquad \qquad \qquad \qquad \qquad \qquad \qquad \\ \mathrm{E} _ {1} ^ {*} \otimes \mathrm{F} _ {1} ^ {*} \xrightarrow {\mu} (\mathrm{E} _ {1} \otimes \mathrm{F} _ {1}) ^ {*} \end{array}\tag{25}
\]

is commutative.

COROLLARY 1. If one of the modules E, F is projective and finitely generated, the canonical homomorphism (23) is bijective.

COROLLARY 2. Let \(\mathbf{E}_1, \mathbf{E}_2\) be two finitely generated projective C-modules, \(u_1\) an endomorphism of \(\mathbf{E}_1\) and \(u_2\) an endomorphism of \(\mathbf{E}_2\) ; then

\[
\operatorname{Tr} (u _ {1} \otimes u _ {2}) = \operatorname{Tr} (u _ {1}) \operatorname{Tr} (u _ {2}).\tag{26}
\]

By linearity, it suffices to consider the case where \(u_{1}\) is of the form \(x_{1} \mapsto \langle x_{1}, x_{1}^{*} \rangle y_{1}\) and \(u_{2}\) of the form \(x_{2} \mapsto \langle x_{2}, x_{2}^{*} \rangle y_{2}\) ; then the image of \(x_{1} \otimes x_{2}\) under \(u_{1} \otimes u_{2}\) is by definition

\[
\langle x _ {1}, x _ {1} ^ {*} \rangle \langle x _ {2}, x _ {2} ^ {*} \rangle (y _ {1} \otimes y _ {2}) = \langle x _ {1} \otimes x _ {2}, x _ {1} ^ {*} \otimes x _ {2} ^ {*} \rangle (y _ {1} \otimes y _ {2})
\]

\(x_{1}^{*}\otimes x_{2}^{*}\) being canonically identified under \(\mu\) with an element of \((\mathrm{E}_{1}\otimes\mathrm{E}_{2})^{*}\) . As \(\langle y_{1}\otimes y_{2},x_{1}^{*}\otimes x_{2}^{*}\rangle=\langle y_{1},x_{1}^{*}\rangle\langle y_{2},x_{2}^{*}\rangle\) , formula (26) follows in this case from (17).

Remark. If E, F, G are any three C-modules, it is immediately verified that the diagram

\[
\begin{array}{c} \mathrm{E} ^ {*} \otimes \mathrm{F} ^ {*} \otimes \mathrm{G} ^ {*} \xrightarrow {\mu \otimes 1} (\mathrm{E} \otimes \mathrm{F}) ^ {*} \otimes \mathrm{G} ^ {*} \\ 1 \otimes \mu \Biggl \downarrow \\ \mathrm{E} ^ {*} \otimes (\mathrm{F} \otimes \mathrm{G}) ^ {*} \xrightarrow [ \mu ]{} (\mathrm{E} \otimes \mathrm{F} \otimes \mathrm{G}) ^ {*} \end{array}\tag{27}
\]

is commutative, by virtue of formula (24).

We note also that, without any hypothesis on the C-modules E, F, there are canonical isomorphisms

\begin{enumerate}
\def\labelenumi{(\arabic{enumi})}
\setcounter{enumi}{27}
\tightlist
\item
\end{enumerate}

\[
(\mathbf {E} \otimes \mathbf {F}) ^ {*} \rightarrow \operatorname{Hom} (\mathbf {E}, \mathbf {F} ^ {*})\tag{29}
\]

\[
(\mathbf {E} \otimes \mathbf {F}) ^ {*} \rightarrow \operatorname{Hom} (\mathbf {F}, \mathbf {E} ^ {*})
\]

which are just the isomorphism (6) and (5) of no. 1 for \(\mathbf{G} = \mathbf{C}\) , \(\mathbf{A} = \mathbf{B} = \mathbf{C}\) .

Thus a canonical one-to-one correspondence has been defined between the bilinear forms on \(E \times F\) , the homomorphisms of E into \(F^{*}\) and the homomorphisms of F into \(E^{*}\) : if u (resp. v) is a homomorphism of E into \(F^{*}\) (resp. of F into \(E^{*}\) ), the corresponding bilinear form is given by

\[
(x, y) \mapsto \langle y, u (x) \rangle \quad (\text { resp. } (x, y) \mapsto \langle x, v (y) \rangle).
\]

